\begin{encadre}[colframe=orange!70]
\textbf{Ce que change le niveau 2 (frottement, entaille des rainures, granulats et mortier).}
\begin{enumerate}[nosep]
\item \textbf{Les conclusions du niveau 1 tiennent} : amplification près de la surface, effet nul
      au-delà de 20 à \SI{40}{mm}, résultats du TFE en profondeur inchangés.
\item \textbf{Mais le niveau 1 sous-estimait la sévérité superficielle}, surtout pour les
      rainures. L'entaille multiplie encore $\epsilon_1$ par 2,1 à \SI{1}{mm}
      (arêtes vives ; 1,6 avec des arêtes arrondies), avec deux points chauds :
      le bord supérieur des plateaux et le fond de la rainure. Avec un freinage courant ($\mu=0{,}3$), une surface rainurée
      atteint 3,7 à
      4,0 fois la déformation
      de la surface lisse à \SI{1}{mm}, contre 1,7 au niveau 1.
\item \textbf{Le frottement} augmente $\epsilon_1$ superficielle de 46\,\%
      (surface lisse) à 90\,\% (MPD \SI{1.5}{mm})
      pour $\mu=0{,}6$ : il accentue l'effet de la texture.
\item \textbf{Dans le mortier} entre granulats, la déformation locale sous une texture vaut
      2,7 à 3,7 fois la déformation moyenne de l'enrobé.
      Cette concentration existe aussi, et plus forte, sous une surface lisse : c'est une propriété
      du matériau. L'écart \emph{relatif} entre texture et surface lisse est donc atténué dans le
      mortier (facteur 0,65 à 0,77),
      alors que la déformation \emph{absolue} y est bien plus forte.
\item \textbf{Confiance} : les sens de variation sont sûrs ; les facteurs du niveau 2 sont
      semi-quantitatifs (facteur 1,5 à 2), car ils dépendent de paramètres à mesurer : distance
      critique, arrondi des arêtes, contraste granulat/mortier (\S\ref{sec:n2_confiance}).
\end{enumerate}
\end{encadre}
